\documentclass[conference]{IEEEtran}
\usepackage[utf8]{inputenc}
\usepackage[T1]{fontenc}
\usepackage[brazilian]{babel}
\usepackage{amsmath}
\usepackage{graphicx}
\usepackage{float}
\usepackage{booktabs}
\usepackage[hidelinks]{hyperref}

\begin{document}

\title{Sistema de Assinatura Digital com RSA-OAEP e RSA-PSS\\
\large Trabalho de Implementação 2 --- CIC0201 Segurança Computacional}

\author{
\IEEEauthorblockN{Ana Luísa Reis Nascente}
\IEEEauthorblockA{Matrícula: 211045688}
\and
\IEEEauthorblockN{Gabriel de Sousa}
\IEEEauthorblockA{Matrícula: 211056000}
\and
\IEEEauthorblockN{Marina Pimentel Moreno}
\IEEEauthorblockA{Matrícula: 222014071}
\and
\IEEEauthorblockN{Pedro de Paula Campos}
\IEEEauthorblockA{Matrícula: 231036050}
}

\maketitle

\begin{abstract}
% TODO (até 5 linhas): o que foi implementado (chaves RSA-2048 com
% Miller-Rabin, OAEP e PSS com SHA3-256/MGF1), como foi validado (testes de
% adulteração e interoperabilidade) e principal resultado.
\end{abstract}

\begin{IEEEkeywords}
RSA, OAEP, PSS, assinatura digital, SHA-3
\end{IEEEkeywords}

\section{Introdução}
% Contexto: criptografia assimétrica na disciplina; diferença entre usar o
% RSA para cifragem e para assinatura \cite{roteiro,rfc8017}.
% Última frase: organização do relatório (Seções II, III, IV).

\section{Fundamentação Teórica}
\label{sec:fundamentacao}

\subsection{RSA e geração de chaves}
% n = pq, lambda(n), e = 65537, d, CRT; Miller-Rabin \cite{sp80056b,fips1865}.

\subsection{Função hash SHA3-256 e MGF1}
% Esponja Keccak \cite{fips202}; MGF1 como gerador de máscara.

\subsection{RSA-OAEP}
% Codificação EME-OAEP, Feistel de duas rodadas, IND-CCA2 \cite{rfc8017}.

\subsection{RSA-PSS}
% EMSA-PSS, salt, trailer 0xbc; por que não é ``cifrar o hash''.

\subsection{Análise de segurança}
% Parte V: textbook RSA (determinismo, maleabilidade, e pequeno),
% papel do OAEP e do PSS, comparação com Ed25519 \cite{rfc8032}.
% Ver parte5_analise/README.md.

\section{Ambiente Experimental e Análise de Resultados}
\label{sec:experimental}

\subsection{Descrição do Cenário}
% Hardware, SO, versão do Python, bibliotecas permitidas (hashlib, secrets),
% arquitetura do código (partes I--IV), formato das chaves e da assinatura.

\subsection{Análise de Resultados}
% Parte I: tempo de geração das chaves, rodadas de Miller-Rabin.
% Parte II: ida e volta OAEP, tamanho máximo (190 bytes), detecção de
%           ciphertext adulterado e rótulo errado.
% Parte III: assinaturas probabilísticas, Base64.
% Parte IV: tabela com os casos (a) byte do arquivo, (b) byte da assinatura,
%           (c) chave pública -> resultado esperado x obtido.
% Interoperabilidade com a biblioteca cryptography.
% Problemas encontrados e soluções.

\section{Conclusões}
\label{sec:conclusoes}
% Entre 5 e 10 linhas.

\begin{thebibliography}{9}

\bibitem{roteiro}
P.~S.~Barreto, ``Trabalho de Implementação 2 -- Sistema de Assinatura Digital e Verificação Segura de Arquivos,'' CIC0201 Segurança Computacional, Universidade de Brasília, 2026.

\bibitem{rfc8017}
K.~Moriarty, B.~Kaliski, J.~Jonsson e A.~Rusch, ``PKCS \#1: RSA Cryptography Specifications Version 2.2,'' RFC 8017, IETF, 2016.

\bibitem{fips202}
NIST, ``SHA-3 Standard: Permutation-Based Hash and Extendable-Output Functions,'' FIPS PUB 202, 2015.

\bibitem{fips1865}
NIST, ``Digital Signature Standard (DSS),'' FIPS PUB 186-5, 2023.

\bibitem{sp80056b}
NIST, ``Recommendation for Pair-Wise Key-Establishment Using Integer Factorization Cryptography,'' SP 800-56B Rev.~2, 2019.

\bibitem{rfc8032}
S.~Josefsson e I.~Liusvaara, ``Edwards-Curve Digital Signature Algorithm (EdDSA),'' RFC 8032, IETF, 2017.

\end{thebibliography}

\end{document}
